% Source: https://artofproblemsolving.com/community/c3h1533045_define_n
% Stable identity: aops-topic-1533045-q01
Clearing denominators,
\[
23(7240+a+n)=67(5043+100b-n)\iff 90n=67\cdot5043-23\cdot7240+6700b-23a=171361+6700b-23a .
\]
\textbf{Units digit.} The left side is a multiple of $10$, so $171361-23a\equiv1-3a\equiv0\pmod{10}$, i.e.\ $3a\equiv1\pmod{10}$, which gives $a=7$. Then
\[
90n=171200+6700b\iff 9n=17120+670b .
\]
\textbf{Divisibility by $9$.} Modulo $9$, $17120\equiv1+7+1+2+0\equiv2$ and $670\equiv4$, so $2+4b\equiv0\pmod9$, i.e.\ $4b\equiv7\pmod9$. Multiplying by $7$ (as $4\cdot7\equiv1$): $b\equiv49\equiv4\pmod9$, so $b=4$. Then $9n=17120+2680=19800$ and $n=\boxed{2200}$.

Check: $\dfrac{7247+2200}{5443-2200}=\dfrac{9447}{3243}=\dfrac{141\cdot67}{141\cdot23}=\dfrac{67}{23}$.
